
\subsection{滤过 $A_{n,K}$ 结构}
\label{subsec:7-2-6}
\liftIndexRangeStart{I248}%
\liftIndexRangeStart{I249}%

在 7.2.6--7.2.8 小节中，
我们处理第~\ref{subsec:7-2-3} 小节提出的问题。
本小节把 $A_K$ 结构的概念推广到滤过情形。
这一构造仍采用与第 7.2.4、7.2.5 小节相似的归纳方法。
设 $\mathfrak G$ 为幺半群，并设
\[
\pi_{\mathfrak G,G}\colon
\mathfrak G\longrightarrow
G\subset\mathbb R_{\ge0}\times2\mathbb Z
\]
为到满足条件 3.1.6 的子幺半群 $G$ 的同态，
使 $G$ 中每个元素的逆像是有限集，
且 $\pi_{\mathfrak G,G}^{-1}(\{(0,0)\})$
仅含一个元素 $\beta_0$。

\begin{definition}{7.2.61}\label{def:7-2-61}
对 $\beta\in\mathfrak G$ 且 $\beta\ne\beta_0$，定义
\[
\|\beta\|=
\sup\left\{n\ \middle|\
\begin{gathered}
\exists\,\beta_i\in\mathfrak G\setminus\{\beta_0\},\\[-2pt]
\sum_{i=1}^{n}\beta_i=\beta
\end{gathered}\right\}
+[\pi_1(\beta)]-1,
\]
并置 $\|\beta_0\|=-1$。
这里 $\pi_1(\beta)$ 是
$\pi_{\mathfrak G,G}(\beta)$ 的 $\mathbb R_{\ge0}$ 分量，
$[x]$ 表示不超过 $x$ 的最大整数。
\end{definition}

注意有
\begin{equation}
\|\beta_1+\beta_2\|
\ge\|\beta_1\|+\|\beta_2\|+1.
\tag{7.2.62}\label{eq:7-2-62}
\end{equation}

\begin{definition}{7.2.63}\label{def:7-2-63}
在
$(\mathfrak G\times\mathbb Z_{\ge0})
\setminus\{(\beta_0,0)\}$
上定义偏序 $<$：当且仅当以下两种情形之一成立时，
记 $(\beta_1,k_1)>(\beta_2,k_2)$：

\noindent\textup{(7.2.64.1)}\quad
$\|\beta_1\|+k_1>\|\beta_2\|+k_2$；

\noindent\textup{(7.2.64.2)}\quad
$\|\beta_1\|+k_1=\|\beta_2\|+k_2$，
且 $\|\beta_1\|>\|\beta_2\|$。

若
$\|\beta_1\|+k_1=\|\beta_2\|+k_2$
且 $\|\beta_1\|=\|\beta_2\|$，
记 $(\beta_1,k_1)\sim(\beta_2,k_2)$。
若 $(\beta_1,k_1)>(\beta_2,k_2)$
或 $(\beta_1,k_1)\sim(\beta_2,k_2)$，
记 $(\beta_1,k_1)\gtrsim(\beta_2,k_2)$。

对非负整数 $n,n',m$，
记号 $(\beta,k)<(n,k)$、
$(m,k)<(n',k')$ 也按类似方式使用；
当比较的是两个整数对 $(m,k)$ 与 $(n',k')$ 时，
$\sim$ 就是通常意义下的相等。
\end{definition}

下列引理由定义立即得到。
\begin{lemma}{7.2.65}\label{lem:7-2-65}
若
$(\beta_i,k_i)\in
(\mathfrak G\times\mathbb Z_{\ge0})
\setminus\{(\beta_0,0)\}$（$i=1,2$），
且 $k_1+k_2\ge1$，则
\[
(\beta_1,k_1)\lesssim
(\beta_1+\beta_2,k_1+k_2-1).
\]
若其中成立等价关系 $\sim$，
则 $(\beta_2,k_2)=(\beta_0,1)$。%
\translatornote{原书引理未写 $k_1+k_2\ge1$；
若两者均为零，比较式右端的第二坐标是 $-1$，
不属于刚定义的偏序的定义域。此处只补明有意义的范围；
后文实际使用的是 $k_1+k_2=k+1\ge1$。
参见疑误 A-0438。}
\end{lemma}

现在定义 $\mathfrak G$-滤过 $A_{n,K}$ 结构。
设 $C$ 为自由滤过分次 $\Lambda_{0,\mathrm{nov}}$-模。
取分次 $\mathbb Q$-向量空间 $\overline C$，使
$\overline C\widehat\otimes\Lambda_{0,\mathrm{nov}}\cong C$；
这里 $\widehat\otimes$ 表示代数张量积的完备化。
假设 $\overline C$ 是上链复形，
其上边界算子记为 $\delta$。
设对每个
$(\beta,k)\in
(\mathfrak G\times\mathbb Z_{\ge0})
\setminus\{(\beta_0,0)\}$，
都有次数为 $+1$ 的运算
\begin{equation}
\mathfrak m_{k,\beta}\colon
B_k\overline C[1]\longrightarrow\overline C[1].
\tag{7.2.66}\label{eq:7-2-66}
\end{equation}

\begin{definition}{7.2.67}\label{def:7-2-67}
设 $n,K$ 为非负整数，$\mathfrak G$ 如上。
$C$ 上的一个 $\mathfrak G$-滤过 $A_{n,K}$ 代数结构
\liftIndex{I245}%
是对满足
$(\beta,k)\lesssim(n,K)$、
$(\beta,k)\ne(\beta_0,0)$ 的指标给出运算
$\mathfrak m_{k,\beta}$，
并要求对每个这样的 $(\beta,k)$，有
\begin{equation}
\sum_{\substack{\beta_1+\beta_2=\beta,\ \beta_i\in\mathfrak G\\
                  k_1+k_2=k+1}}
\sum_i(-1)^{\deg'\mathbf x_i^{(1)}}
\mathfrak m_{k_2,\beta_2}
\bigl(\mathbf x_i^{(1)},
\mathfrak m_{k_1,\beta_1}(\mathbf x_i^{(2)}),
\mathbf x_i^{(3)}\bigr)=0,
\tag{7.2.68}\label{eq:7-2-68}
\end{equation}
其中
$\Delta^2(\mathbf x)=
\sum_i\mathbf x_i^{(1)}\otimes
\mathbf x_i^{(2)}\otimes\mathbf x_i^{(3)}$。
\end{definition}
\liftIndexRangeEnd{I249}%

由条件 3.1.6，式~\eqref{eq:7-2-68} 左端是有限和；
由引理~\ref{lem:7-2-65}，它是良定义的。
以下固定 $\mathfrak G$，
因而简称 $\mathfrak G$-滤过 $A_{n,K}$ 结构为
滤过 $A_{n,K}$ 结构，其他相关名称亦然。
下面的引理由定义容易得出。

\begin{lemma}{7.2.69}\label{lem:7-2-69}
设 $C$ 对每对非负整数 $n,K$ 都有一个
滤过 $A_{n,K}$ 代数结构，
且每当 $(n,K)<(n',K')$ 时，
滤过 $A_{n',K'}$ 结构限制后
与滤过 $A_{n,K}$ 结构相同。
则
\[
\mathfrak m_k=
\sum_{\beta\in\mathfrak G}
T^{\lambda(\beta)}e^{\mu(\beta)/2}
\mathfrak m_{k,\beta}
\]
是一个滤过 $A_\infty$ 结构。
这里
$\lambda(\beta)=\pi_{\mathfrak G,G,1}(\beta)
\in\mathbb R_{\ge0}$
与
$\mu(\beta)=\pi_{\mathfrak G,G,2}(\beta)
\in2\mathbb Z$
分别是 $\pi_{\mathfrak G,G}(\beta)$ 的第一、第二分量。
\end{lemma}

设 $C,C'$ 为滤过 $A_{n,K}$ 代数，
并考虑一族次数为零的 $\mathbb Q$-线性同态
\[
\mathfrak f_{k,\beta}\colon
B_k\overline C[1]\longrightarrow\overline C'[1].
\]
假设 $\mathfrak f_{0,\beta_0}=0$。
\begin{definition}{7.2.70}\label{def:7-2-70}
若对每个 $(\beta,k)\lesssim(n,K)$ 都有下式，
则称 $\{\mathfrak f_{k,\beta}\}$
定义一个滤过 $A_{n,K}$ 同态\liftIndex{I246}\liftIndex{I250}：
\begin{equation*}
\begin{aligned}
&\sum_{\substack{m,i,\ \beta'+\beta_1+\cdots+\beta_m=\beta\\
                   k_1+\cdots+k_m=k}}
\mathfrak m_{m,\beta'}
\bigl(\mathfrak f_{k_1,\beta_1}(\mathbf x_i^{(1)}),
\ldots,
\mathfrak f_{k_m,\beta_m}(\mathbf x_i^{(m)})\bigr)\\
&\quad=
\sum_{\substack{\beta_1+\beta_2=\beta\\
                   k_1+k_2=k+1}}
\sum_i(-1)^{\deg'\mathbf x_i^{(1)}}
\mathfrak f_{k_2,\beta_2}
\bigl(\mathbf x_i^{(1)},
\mathfrak m_{k_1,\beta_1}(\mathbf x_i^{(2)}),
\mathbf x_i^{(3)}\bigr).
\end{aligned}
\end{equation*}
\end{definition}

对于滤过 $A_{n,K}$ 同态，
也可证明与引理~\ref{lem:7-2-69} 类似的结果。

还可按显然方式定义滤过 $A_{n,K}$ 代数 $C$ 的
$[0,1]\times C$ 模型。
若 $C$ 是滤过 $A_{n,K}$ 代数，
则第~\ref{sec:4-2} 节引入的
$\operatorname{Poly}([0,1],C)$ 与 $C^{[0,1]}$
都是滤过 $A_{n,K}$ 代数，
也都是 $[0,1]\times C$ 的模型。
可借此定义滤过 $A_{n,K}$ 同态之间的
滤过 $A_{n,K}$ 同伦\liftIndex{I247}\liftIndex{I251}。
利用与定理 4.2.34 类似的结果，
可证明它是等价关系。
从而可定义滤过 $A_{n,K}$ 同伦等价，
定理 4.2.45 的类似结论也成立。
\liftIndexRangeEnd{I248}%

\begin{remark}{7.2.71}\label{rem:7-2-71}
证明定理 4.2.34 与 4.2.45 的 $A_{n,K}$ 版本时，
我们按偏序 $<$ 归纳。
这与第 4 章滤过 $A_\infty$ 版本的证明略有不同：
那里采用了双重归纳，
分别沿数目滤过与能量滤过进行。
不过，滤过 $A_\infty$ 版本的上述两个定理，
也可以一次性地按偏序 $<$ 归纳证明。
\end{remark}

现在证明下列结果。
\begin{theorem}{7.2.72}\label{thm:7-2-72}
设 $C_1$ 为滤过 $A_{n,K}$ 代数，
$C_2$ 为滤过 $A_{n',K'}$ 代数，
且 $(n,K)<(n',K')$。
设 $\mathfrak h\colon C_1\to C_2$
为滤过 $A_{n,K}$ 同态，
并假设它是滤过 $A_{n,K}$ 同伦等价。
则可在 $C_1$ 上赋予一个
扩张既有 $A_{n,K}$ 代数结构的
滤过 $A_{n',K'}$ 代数结构，
且可把 $\mathfrak h$ 扩张成
$C_1\to C_2$ 的滤过 $A_{n',K'}$ 同伦等价。
\end{theorem}

\begin{proof}
为证明定理，需要研究扩张滤过 $A_{n,K}$ 结构与同态的障碍。
由偏序 $<$ 的定义，只需处理以下两种情形：

\noindent\textup{(7.2.73.1)}\quad
$(n,K)<(n',K')=(n+1,K-1)$；

\noindent\textup{(7.2.73.2)}\quad
$(n,K)=(n,0)<(n',K')=(0,n+1)$。

回忆第 4.5.1 小节开头定义的上边界算子
\[
\delta_1(\varphi)
=\overline{\mathfrak m}_1\circ\varphi
+(-1)^{\deg\varphi+1}
\varphi\circ\widehat{\overline{\mathfrak m}}_1,
\]
其中
$\widehat{\overline{\mathfrak m}}_1\colon
B\overline C[1]\to B\overline C[1]$
是如式 (3.2.1) 那样由
$\overline C$ 上的 $\overline{\mathfrak m}_1$
诱导的余导子。

\begin{lemma}{7.2.74}\label{lem:7-2-74}
设 $(n',K')$ 取
(7.2.73.1) 或 (7.2.73.2) 中的一种情形，
且 $C$ 为滤过 $A_{n,K}$ 代数。
则对每个满足 $\|\beta\|=n'$ 的 $\beta$，
存在障碍类\liftIndex{I252}
\[
[o_{K',\beta}(C)]\in
H\bigl(\operatorname{Hom}
(B_{K'}\overline C[1],\overline C[1]),
\delta_1\bigr),
\]
并且所有这些类均为零，
当且仅当存在扩张既有滤过 $A_{n,K}$ 结构的
滤过 $A_{n',K'}$ 结构。
此外，若 $C\to C'$ 是滤过 $A_{n,K}$ 同伦等价，
则这一同伦等价诱导的同构
\[
\begin{aligned}
&H\bigl(\operatorname{Hom}
(B_{K'}\overline C[1],\overline C[1]),\delta_1\bigr)\\
&\qquad\cong
H\bigl(\operatorname{Hom}
(B_{K'}\overline C'[1],\overline C'[1]),\delta_1\bigr)
\end{aligned}
\]
把 $[o_{K',\beta}(C)]$ 映到
$[o_{K',\beta}(C')]$。
\end{lemma}

\begin{proof}
令 $(\|\beta\|,k)=(n',K')$。
考虑
\begin{equation}
\sum_{\substack{\beta_1+\beta_2=\beta,\ k_1+k_2=k+1\\
 (k_i,\beta_i)\ne(k,\beta)}}
\sum_i(-1)^{\deg'\mathbf x_i^{(1)}}
\mathfrak m_{k_2,\beta_2}
\bigl(\mathbf x_i^{(1)},
\mathfrak m_{k_1,\beta_1}(\mathbf x_i^{(2)}),
\mathbf x_i^{(3)}\bigr),
\tag{7.2.75}\label{eq:7-2-75}
\end{equation}
其中 $\mathbf x\in B_{K'}\overline C$。
它定义了
\[
o_{K',\beta}(C)\in
\operatorname{Hom}
(B_{K'}\overline C[1],\overline C[1]).
\]
由引理~\ref{lem:7-2-65}，
式~\eqref{eq:7-2-75} 良定义。
与定理 4.5.1 的证明类似，
可知它是 $\delta_1$-闭上链，
故确定上同调类 $[o_{K',\beta}(C)]$。
若所有这些类为零，
便可找到 $\mathfrak m_{K',\beta}$，使
\begin{equation}
\mathfrak m_{1,\beta_0}\circ
\mathfrak m_{K',\beta}
+\mathfrak m_{K',\beta}\circ
\widehat{\mathfrak m}_{1,\beta_0}
+o_{K',\beta}(C)=0.
\tag{7.2.76.1}\label{eq:7-2-76-1}
\end{equation}
回忆 $\overline{\mathfrak m}_1=\mathfrak m_{1,\beta_0}$。
于是可用这些 $\mathfrak m_{K',\beta}$
将原结构扩张成滤过 $A_{n',K'}$ 结构。
反向论证相同；自然性容易验证，留给读者。
\end{proof}

现在完成定理~\ref{thm:7-2-72} 的证明。
设成立的是 (7.2.73.1) 或 (7.2.73.2)。
引理~\ref{lem:7-2-74} 表明，
$C_1$ 上的滤过 $A_{n,K}$ 结构
可扩张为滤过 $A_{n',K'}$ 结构。
还须扩张同态 $\mathfrak h\colon C_1\to C_2$，
即对每个 $\|\beta\|=n'$ 定义
\[
\mathfrak h_{K',\beta}\colon
B_{K'}\overline C_1[1]
\longrightarrow\overline C_2[1].
\]
这些映射构成滤过 $A_{n',K'}$ 同态的条件等价于
\begin{equation}
\mathfrak m_{1,\beta_0}\circ
\mathfrak h_{K',\beta}
+\mathfrak h_{K',\beta}\circ
\widehat{\mathfrak m}_{1,\beta_0}
+o_{K',\beta}(\mathfrak h)=0.
\tag{7.2.76.2}\label{eq:7-2-76-2}
\end{equation}
这里 $o_{K',\beta}(\mathfrak h)$
按与第 4.5 节及引理~\ref{lem:7-2-74} 的证明类似的方式定义。
可证明它是闭上链，
因而定义类\liftIndex{I253}
\[
[o_{K',\beta}(\mathfrak h)]\in
H\bigl(\operatorname{Hom}
(B_{K'}\overline C_1[1],\overline C_2[1])\bigr).
\]
要证明可选择 $C_1$ 上的运算
$\mathfrak m_{K',\beta}$，
使 $[o_{K',\beta}(\mathfrak h)]=0$。
注意 $o_{K',\beta}(\mathfrak h)$ 中
涉及这一运算的项是
\[
\mathfrak h_{1,\beta_0}\circ
\mathfrak m_{K',\beta}.
\]
由引理~\ref{lem:7-2-74} 的证明，
可以给 $\mathfrak m_{K',\beta}$
加上复形
$\operatorname{Hom}
(B_{K'}\overline C_1[1],\overline C_1[1])$
中的任意闭上链，
所得运算仍满足式~\eqref{eq:7-2-76-1}。
另一方面，由同伦等价假设，
$\mathfrak h_{1,\beta_0}\circ$
诱导上同调同构
\[
\begin{aligned}
&H\bigl(\operatorname{Hom}
(B_{K'}\overline C_1[1],\overline C_1[1])\bigr)\\
&\qquad\longrightarrow
H\bigl(\operatorname{Hom}
(B_{K'}\overline C_1[1],\overline C_2[1])\bigr).
\end{aligned}
\]
因此可选择 $\mathfrak m_{K',\beta}$，
使 $[o_{K',\beta}(\mathfrak h)]=0$，
进而把 $\mathfrak h$ 扩张为滤过 $A_{n',K'}$ 同态。
由 $A_{n',K'}$ 版本的 Whitehead 定理，
这一同态还是滤过 $A_{n',K'}$ 同伦等价；
该版本可仿照第 4.5 节证明。
定理~\ref{thm:7-2-72} 证毕。
\end{proof}
